\documentclass[11pt]{report}
\usepackage[margin=0.59in]{geometry}
\usepackage{graphicx}
\usepackage{tikz}
\usetikzlibrary{calc}
\usepackage{hyperref}

\definecolor{colorrds}{HTML}{0060A1} % corrected the RGB values to represent royal purple more accurately
\newcommand{\mediumsize}{\fontsize{15pt}{16pt}\selectfont}

\begin{document}
\begin{titlepage}

    \begin{tikzpicture}[remember picture,overlay]
        \draw[line width=10 pt, color=colorrds]
            ($(current page.north west) + (0.30in,-0.30in)$) rectangle
            ($(current page.south east) + (-0.30in,0.30in)$);
    \end{tikzpicture}

%     \begin{minipage}[t]{0.95\textwidth}
%         \hfill
%         \raggedleft
%         \textbf{ Jane Doe} \\
%         New York, NY \\
%         \today\\
%         \href{mailto:jane.doe@email.com}{jane.doe@email.com} \\
%         +1 436-823-9879 \\
%         \url{http://jane-doe.com} \\
%         \url{http://linkedin.com/in/jane-doe}
%     \end{minipage}

% \raggedright \textbf{Maria Winter, Ph.D.} \\ Department of Political Science \\ Harvard University \\ Cambridge, MA 02138, USA \\ \href{mailto:maria.winter@harvard.edu}{maria.winter@harvard.edu}\\

\vspace{0.7em}

\raggedright Estimados alumnos de segundo grado:\\

Les doy la más cálida y profunda bienvenida a este ciclo escolar 2026-2027 en la escuela secundaria Rafael Díaz Serdán. Mi nombre es Julio, y este año tendré nuevamente el privilegio de guiarlos a través del vasto y fascinante universo de las Matemáticas. Si durante su primer año comenzamos a quitarle el velo a los números para descubrir que no son herramientas de tortura, sino el lenguaje en el que está escrito el cosmos, este segundo año nos adentraremos en aguas mucho más profundas. Ha llegado el momento de dar un salto hacia la abstracción, de abandonar la orilla de lo conocido y atrevernos a navegar en el océano del álgebra y la geometría, donde las ideas puras toman forma y nos revelan los secretos más íntimos de la realidad.

\vspace{0.7em}

A lo largo de los próximos meses, notarán que nuestro lenguaje comenzará a transformarse. Empezaremos a invitar a las letras a bailar con los números. Para muchos, este es el momento en que las Matemáticas parecen volverse un enigma indescifrable, pero les pido que confíen en mí y abran su mente. Una "x" o una "y" en una ecuación no están ahí para confundirlos; son, por el contrario, contenedores de posibilidades infinitas, son misterios esperando ser resueltos, son el símbolo más puro de la curiosidad humana. El álgebra no es más que el arte de encontrar lo que está oculto, de equilibrar la balanza de la verdad y de trazar un camino lógico desde la oscuridad de la incertidumbre hasta la luz de la certeza. Cuando resolvemos una ecuación compleja, no estamos simplemente llenando una hoja con garabatos; estamos entrenando a nuestra mente para enfrentar el caos del mundo y encontrar en él un orden perfecto, una estructura, una solución.

\vspace{0.7em}

Nuestra aula seguirá siendo un laboratorio de pensamiento vivo. Aquí no hay espacio para la repetición autómata ni para la memorización carente de sentido. Exijo de ustedes algo mucho más valioso: su capacidad de asombro y su tenacidad. Las Matemáticas de segundo grado les exigirán paciencia, les pedirán que se detengan a observar un problema desde distintos ángulos antes de atacarlo, y les enseñarán que el camino más directo no siempre es el más evidente. Aprenderemos a leer el espacio a través de la geometría, comprendiendo que las formas que construimos y habitamos no son accidentes, sino respuestas calculadas a las fuerzas de la naturaleza. Entenderemos que detrás de la arquitectura de nuestra escuela, del arco de un puente o de la trayectoria de un balón en el aire, existe un esqueleto invisible de triángulos, ángulos y proporciones que sostienen la existencia misma.

\vspace{0.7em}

Les invito a entrar a este salón despojados del miedo a equivocarse. El error es el cincel con el que esculpimos el conocimiento; sin él, la mente se estanca en la comodidad. Atrévanse a proponer soluciones descabelladas, a debatir con la lógica en la mano, a cuestionar cada paso y a no descansar hasta que la mente experimente esa chispa irreemplazable de comprensión pura. Este viaje que hoy emprendemos no solo busca prepararlos para un examen, sino forjar en ustedes un intelecto afilado, crítico y libre, capaz de descifrar cualquier enigma que la vida les presente.

\vspace{0.7em}

Con profunda emoción por los descubrimientos que nos aguardan, les doy la bienvenida al arte del pensamiento abstracto. Bienvenidos a Matemáticas de segundo grado.

\vspace{0.7em}

\raggedright Atentamente,\\
\textbf{Julio Melchor}

\end{titlepage}
\end{document}
